% Generated draft: requires mathematical and source review.
\documentclass[11pt]{article}
\usepackage[margin=25mm]{geometry}
\usepackage{fontspec}
\setmainfont{Latin Modern Roman}
\usepackage{amsmath,amssymb,mathtools,unicode-math}
\setmathfont{Latin Modern Math}
\usepackage{enumitem,needspace,xurl,hyperref}
\setlength{\parindent}{0pt}
\setlength{\parskip}{5pt}
\setlength{\emergencystretch}{3em}
\newcommand{\N}{\mathbb{N}}
\newcommand{\Z}{\mathbb{Z}}
\newcommand{\Q}{\mathbb{Q}}
\newcommand{\R}{\mathbb{R}}
\newcommand{\C}{\mathbb{C}}
\newcommand{\F}{\mathbb{F}}
\newcommand{\A}{\mathbb{A}}
\newcommand{\PP}{\mathbb{P}}
\newcommand{\E}{\mathbb{E}}
\newcommand{\eps}{\varepsilon}
\newcommand{\dd}{\,\mathrm{d}}
\begin{document}
\section*{UnsolvedMath ID: 20000388}
\textbf{Source problem code:} catalogue entry
\subsection{Definitions and mathematical statement}
The source proposes studying complete intersections and examples in which the target space varies. A complete intersection in \(\PP^n\) is a closed subscheme cut out by a regular sequence of homogeneous equations of the expected codimension; a smooth complete-intersection curve cut out by \(n-1\) hypersurfaces has dimension one. The source sentence does not specify a base field, distribution on equations, statistic, limit, or meaning of ``moving target space'', and therefore gives a research direction rather than a unique quantified conjecture.

The supplied research report studies the following narrower point-count setting. Fix a prime power \(q\). For each fixed ambient dimension \(n\), consider smooth curves over \(\F_q\) cut out by \(n-1\) hypersurfaces in \(\PP^n\), with the defining degrees tending to infinity in the Bucur--Kedlaya limiting model referenced in the report. Write \(\mu_{q,n}\) for the limiting mean of \(\#C(\F_q)\), and \(L_{q,n}\) for its reported limiting binomial law. The candidate theorem asserts
\[
 \mu_{q,1}=q+1,\qquad
 \mu_{q,n+1}<\mu_{q,n}\quad(n\ge1),\qquad
 \lim_{n\to\infty}\mu_{q,n}=\lambda_q
 =(q+1)\prod_{j=3}^{\infty}(1-q^{-j}),
\]
and
\[
 \lim_{n\to\infty}d_{\mathrm{TV}}
 \bigl(L_{q,n},\operatorname{Pois}(\lambda_q)\bigr)=0.
\]
Here \(\operatorname{Pois}(\lambda)\) has mass \(e^{-\lambda}\lambda^m/m!\) at \(m\in\Z_{\ge0}\), and total variation on nonnegative integers is \(\tfrac12\sum_{m\ge0}|\mu(\{m\})-\nu(\{m\})|\). The dimension-one endpoint corresponds to \(\PP^1\) with no cutting equations. These are iterated limits: first the hypersurface degrees tend to infinity for each fixed \(n\), and only then does \(n\to\infty\). No simultaneous degree/dimension limit is asserted.

The source record does not provide the explicit fixed-\(n\) binomial parameters or the probability measure and degree-growth hypotheses of the Bucur--Kedlaya model. Accordingly the notation \(\mu_{q,n}\) and \(L_{q,n}\) identifies the quantities in the supplied report, without furnishing a complete independent definition of that model. The report also claims a secondary jet-support theorem giving an exact coefficient-rank obstruction to independent local data for moving-target families; the relevant families, jet orders, linear map and rank formula are absent. That secondary claim cannot be upgraded here to a specific theorem or used to fill in the original meaning of ``moving target''.
\subsection{Short English statement}
Wood asks for complete-intersection and moving-target examples without specifying a precise statistic. The supplied later report studies point counts of smooth complete-intersection curves over a fixed finite field, claiming strictly decreasing limiting means and an iterated Poisson limit as ambient dimension grows. The underlying sampling model and the separate jet-support theorem still require source review.
\subsection{Sources}
\begin{itemize}
\item Local UnsolvedMath registry, permanent ID \(20000388\), source code \texttt{catalogue entry}; AIM algebraic number theory corpus, \texttt{aim-algebraic-number-theory-notes.json}, notes[39]. Attribution: Ulam AI / UnsolvedMath Contributors, CC BY 4.0.
\item AIM workshop \emph{Arithmetic statistics over finite fields and function fields}, Question 10, attributed to Melanie Matchett Wood: \url{https://aimath.org/pastworkshops/arithstatffieldproblems.pdf}. Original remarks and literature fields are empty. The supplied verification identifies this question but provides no further verification details after its introductory colon.
\item The research summary refers to Bucur--Kedlaya complete-intersection limiting laws without a full bibliographic citation or model statement; no missing title, date or formula is invented.
\item Registry dated August 14, 2026: \texttt{partially\_solved}, \texttt{PARTIAL-PROGRESS}; attempt 1 is labelled partial result, rigor proved, novelty confidence low. Its monotonicity, Poisson-limit and jet-support assertions are machine-generated research aids, not peer reviewed, and require independent verification before citation.
\end{itemize}
\end{document}
